% Part: lambda-calculus
% Chapter: lambda-definability
% Section: primitive-recursive-functions

\documentclass[../../../include/open-logic-section]{subfiles}

\begin{document}

\olfileid{lam}{ldf}{prf}
\olsection{ఆదిమ పునరావృత్త ప్రమేయాలు \usetoken{S}{lambda definable}}

ఆదిమ పునరావృత్త ప్రమేయాలు అంటే ప్రాథమిక ప్రమేయాలైన
$\Zero$, $\Succ$, $\Proj{n}{i}$ల నుంచి సంయుక్తం,
ఆదిమ పునరావృత్తి ద్వారా నిర్వచించగల ప్రమేయాలని గుర్తుంచుకోండి.

\begin{lem}
  \ollabel{lem:basic}
  ప్రాథమిక ఆదిమ పునరావృత్త ప్రమేయాలైన $\Zero$, $\Succ$,
  ప్రక్షేప ప్రమేయాలు~$\Proj{n}{i}$
  \tetoken{లాంబ్డాతో నిర్వచించదగినవి}{!!{lambda definable}}.
\end{lem}

\begin{proof}
వాటిని కింది పదాలు \tetoken{లాంబ్డాతో నిర్వచిస్తాయి}{!!{lambda define}}:
\begin{align*}
  \fn{Zero} & \ident \lambd[a][\lambd[fx][x]]\\
  \fn{Succ} & \ident \lambd[a][\lambd[fx][f (a f x)]] \\
  \fn{Proj}^n_i & \ident \lambd[x_0\dots x_{n-1}][x_i]
\end{align*}
\end{proof}

\begin{lem}
  \ollabel{lem:comp} $k$-స్థానిక ప్రమేయం $f$, $n$-స్థానిక
  ప్రమేయాలు $g_0, \dots, g_{k-1}$ వరుసగా $F$, $G_0$,
  \dots, $G_{k-1}$ పదాల ద్వారా
  \tetoken{లాంబ్డాతో నిర్వచించదగినవి}{!!{lambda definable}}
  అనుకుందాం; వాటి సంయుక్తం ద్వారా $h$ నిర్వచితమైందనుకుందాం.
  అప్పుడు $h$ \tetoken{లాంబ్డాతో నిర్వచించదగినది}{!!{lambda definable}}.
  \sourcecorrection{OLTELAMLDFPRF-001}{ఆంగ్ల మూలం $k$ ప్రమేయాలకు ప్రతినిధి పదాలను $G_0,\dots,G_k$గా వ్రాసింది; చివరి సూచిక $k-1$ కావాలి. అదే లెమ్మా ముగింపులో నిర్వచించదగిన ప్రమేయాన్ని $H$గా పేర్కొంది; కింద నిర్మించే ప్రతినిధి పదం $H$, ప్రమేయం $h$ కాబట్టి రెండింటి పాత్రలను సరిచేశాం.}
\end{lem}

\begin{proof}
  $h$ను కింది పదం \tetoken{లాంబ్డాతో నిర్వచిస్తుంది}{!!{lambda define}}:
  \[
  H \ident \lambd[x_0 \dots x_{n-1}][F\, (G_0 x_0 \dots
    x_{n-1}) \dots (G_{k-1} x_0 \dots x_{n-1})]
  \]
  ఇది నిజమని ధృవీకరించడాన్ని అభ్యాసంగా వదిలేస్తున్నాం.
\end{proof}

\begin{prob}
  $H\num{n_0}\dots\num{n_{n-1}} \red \num{h(n_0, \dots,
    n_{n-1})}$ అని చూపి, \olref[lam][ldf][prf]{lem:comp}
  నిరూపణను పూర్తి చేయండి.
\end{prob}

\olref{lem:comp}లో $f$, $g_0$, \dots,~$g_{k-1}$లు
ఆదిమ పునరావృత్త ప్రమేయాలు కావాలని కోరలేదని గమనించండి;
అవి సంపూర్ణమైనవి, \tetoken{లాంబ్డాతో నిర్వచించదగినవి}{!!{lambda definable}}
అయితే చాలు.

\begin{lem}\ollabel{lem:prim}
  $f$ ఒక $n$-స్థానిక ప్రమేయం, $g$ ఒక $n+2$-స్థానిక ప్రమేయం
  అనుకుందాం. అవి వరుసగా $F$, $G$ పదాల ద్వారా
  \tetoken{లాంబ్డాతో నిర్వచించదగినవి}{!!{lambda definable}}
  అనుకుందాం. $f$, $g$ల నుంచి ఆదిమ పునరావృత్తి ద్వారా
  $h$ నిర్వచితమైతే, $h$ కూడా
  \tetoken{లాంబ్డాతో నిర్వచించదగినది}{!!{lambda definable}}.
\end{lem}

\begin{proof}
  $h$ నిర్వచనాన్ని గుర్తుచేసుకుందాం:
  \begin{align*}
    h(x_1, \dots, x_n, 0) &= f(x_1, \dots, x_n)\\
    h(x_1, \dots, x_n, y+1) & = g(x_1, \dots, x_n, y, h(x_1, \dots, x_n, y)).
  \end{align*}
  \sourcecorrection{OLTELAMLDFPRF-002}{ఆంగ్ల మూలం రెండో సమీకరణ కుడి వైపున $h(x_1,\dots,x_n,y,h(x_1,\dots,x_n,y))$ను వ్రాసింది. అది $n+1$-స్థానిక $h$కు $n+2$ ఆర్గ్యుమెంట్లను ఇస్తుంది; దశ ప్రమేయం $g$ కావాలని ముందరి నిర్వచనం, వెంటనే వచ్చే $G$ నిర్మాణం, నిరూపణలోని $g(n,m,h(n,m))$ నిర్ధారిస్తున్నాయి. అందుకే బయట పదాన్ని $g$గా సరిచేశాం.}
  అనౌపచారికంగా, ఈ నిర్వచనం $f(x_1,\dots,x_n)$తో మొదలై,
  ప్రతి దశలో ప్రస్తుత సూచికనూ విలువనూ $g$కు ఇచ్చే నవీకరణను
  $y$ సార్లు పునరావర్తిస్తుంది. ఇది పునరావర్తకంగా పనిచేసే చర్చ్
  సంఖ్యాంకాల నిర్వచనాన్ని గుర్తుచేస్తుంది.

  సరళత కోసం, అదనపు పరామితి ఒక్కటే అయిన సందర్భంలో,
  దాన్ని~$x$గా తీసుకొని నిర్వచనాన్నీ నిరూపణనూ ఇస్తాం.
  $h$ను ఇలా \tetoken{లాంబ్డాతో నిర్వచించవచ్చు}{!!{lambda define}}:
  \begin{align*}
    H \ident & \lambd[x][\lambd[y][\fn{Snd} (y
        D \tuple{\num{0}, F x})]]
    \intertext{ఇక్కడ }
    D \ident & \lambd[p][\tuple{\fn{Succ} (\fn{Fst}\, p), (G x
          (\fn{Fst}\, p) (\fn{Snd}\, p))}]
  \end{align*}
  ప్రతి పునరావర్తన దశలో ఒక క్రమయుగ్మాన్ని స్థితిగా ఉంచుతాం.
  దాని మొదటి అవయవం ప్రస్తుత పునరావర్తన సూచిక~$y$; రెండవ అవయవం
  దానికి సంబంధించిన~$h$ విలువ. ప్రతి దశలో ఈ $y$ సూచికకు,
  $h$కు కొత్త విలువల క్రమయుగ్మాన్ని తయారు చేస్తాం.
  పునరావర్తనాలు పూర్తయ్యాక క్రమయుగ్మపు రెండవ అవయవాన్ని,
  అంటే చివరి $h$ విలువను, తిరిగి ఇస్తాం.
  ఇది ఆదిమ పునరావృత్తిని నిజంగా సూచిస్తుందని ఇప్పుడు నిరూపిస్తాం.

  ఏ $n$, $m$లకైనా $H\,\num{n}\,\num{m} \red
  \num{h(n,m)}$ అని నిరూపించాలి. ముందుగా $D_n \ident
  \Subst{D}{\num{n}}{x}$ అయితే, $D_n^m \tuple{\num{0}, F\, \num{n}} \red
  \tuple{\num{m}, \num{h(n, m)}}$ అని చూపుతాం.
  దీన్ని~$m$పై ఆగమనంతో నిరూపిస్తాం.

  $m=0$ అయితే, $D_n^0 \tuple{\num{0}, F\, \num{n}} \red
  \tuple{\num{0}, \num{h(n, 0)}}$ అని చూపాలి. అయితే
  $D_n^0 \tuple{\num{0}, F\, \num{n}}$ అనేది
  $\tuple{\num{0}, F\, \num{n}}$ మాత్రమే. $F$ అనేది
  $f$ను \tetoken{లాంబ్డాతో నిర్వచిస్తుంది}{!!{lambda define}s} కాబట్టి,
  ఇది $\tuple{\num{0}, \num{f(n)}}$కు తగ్గుతుంది.
  $f(n) = h(n, 0)$ కనుక అది
  $\tuple{\num{0}, \num{h(n,0)}}$ అవుతుంది.

  ఇప్పుడు $D_n^m \tuple{\num{0}, F\, \num{n}} \red
  \tuple{\num{m}, \num{h(n, m)}}$ అని అనుకుందాం.
  అప్పుడు $D_n^{m+1} \tuple{\num{0}, F\, \num{n}} \red
  \tuple{\num{m+1}, \num{h(n, m+1)}}$ అని చూపాలి.
  \begin{align*}
    D_n^{m+1} \tuple{\num{0}, F\, \num{n}}
    & \ident D_n(D_n^{m} \tuple{\num{0}, F\, \num{n}})\\
    & \red D_n\,\tuple{\num{m}, \num{h(n, m)}} \text{\qquad (ఆగమన పరికల్పన వల్ల)}\\
    & \ident (\lambd[p][
      \tuple{
        \fn{Succ} (\fn{Fst}\, p), (G \, \num{n} (\fn{Fst}\, p) (\fn{Snd}\, p))
    }]) \tuple{\num{m}, \num{h(n,m)}}\\
    & \redone
    \tuple{\fn{Succ} (\fn{Fst}\, \tuple{\num{m}, \num{h(n,m)}}),\\
      & \qquad (G \, \num{n} (\fn{Fst}\, \tuple{\num{m}, \num{h(n,m)}}) (\fn{Snd}\, \tuple{\num{m}, \num{h(n,m)}}))}\\
    & \red
    \tuple{\fn{Succ}\, \num{m}, (G \, \num{n} \,\num{m} \, \num{h(n,m)})}\\
    & \red \tuple{\num{m+1}, \num{g(n, m, h(n, m))}}
  \end{align*}
  $g(n, m, h(n, m)) = h(n, m+1)$ కనుక నిరూపణ పూర్తయింది.

  చివరగా, దీన్ని పరిశీలించండి:
  \begin{align*}
    H\,\num n\,\num m &\ident \lambd[x][\lambd[y][\fn{Snd} (y
        (\lambd[p].\tuple{\fn{Succ} (\fn{Fst}\, p), (G\, x\,
          (\fn{Fst}\, p)\, (\fn{Snd}\, p))}) \tuple{\num{0}, F x})]]\\
    & \qquad\,\num n\, \num m\\
    & \red \fn{Snd} (\num m \,
    \underbrace{(\lambd[p].\tuple{\fn{Succ} (\fn{Fst}\, p), (G \,\num n\,
      (\fn{Fst}\, p) (\fn{Snd}\, p))})}_{D_n} \tuple{\num{0}, F \num n})\\
    & \ident \fn{Snd} (\num m \, D_n\, \tuple{\num{0}, F \num n})\\
    & \red \fn{Snd} \,(D_n^m  \tuple{\num{0}, F \num n})\\
    & \red \fn{Snd} \,\tuple{\num{m}, \num{h(n, m)}} \\
    & \red \num{h(n,m)}.
  \end{align*}
\end{proof}

\begin{prop}
  ప్రతి ఆదిమ పునరావృత్త ప్రమేయం
  \tetoken{లాంబ్డాతో నిర్వచించదగినది}{!!{lambda definable}}.
\end{prop}

\begin{proof}
  \olref{lem:basic} ప్రకారం అన్ని ప్రాథమిక ప్రమేయాలు
  \tetoken{లాంబ్డాతో నిర్వచించదగినవి}{!!{lambda definable}};
  \olref{lem:comp}, \olref{lem:prim} ప్రకారం
  \tetoken{లాంబ్డాతో నిర్వచించదగిన}{!!{lambda definable}}
  ప్రమేయాలు సంయుక్తం, ఆదిమ పునరావృత్తి కింద సంవృతమైనవి.
\end{proof}

\end{document}
